% ----- CHAUPAI 11 -----
\subsection*{\deva{एकादश चौपाई} (Eleventh Chaupai)}
\addcontentsline{toc}{subsection}{\deva{एकादश चौपाई} (Eleventh Chaupai) — \deva{लाय सजीवन...}}

\begin{minipage}[t]{0.55\textwidth}
\vspace{0pt}

{\large \linespread{1.5}\selectfont
\deva{लाय सजीवन लखन जियाये।}\\
\deva{श्रीरघुबीर हरषि उर लाये॥}
\par}

\vspace{0.5em}

{\large \linespread{1.5}\selectfont
\telu{లాయ సజీవన లఖన జియాయే।}\\
\telu{శ్రీరఘుబీర హరషి ఉర లాయే॥}
\par}

\vspace{0.5em}

{\large \linespread{1.3}\selectfont
\textit{lāya sajīvana lakhana jiyāye |}\\
\textit{śrīraghubīra haraṣi ura lāye ||}
\par}
\end{minipage}%
\hfill
\begin{minipage}[t]{0.42\textwidth}
\vspace{0pt}
\begin{tcolorbox}[anchorbox]
\textbf{The Ultimate Team Player:} When his teammate fell in battle, Hanuman didn't hesitate; he crossed continents overnight and worked tirelessly to bring the life-saving cure (\textit{Sanjeevani}). He teaches us the absolute importance of resourcefulness, reliability, and showing up unconditionally for our friends in their darkest hours.
\vspace{0.5em}\newline Hanuman worked tirelessly through the night to save his fallen commander, defining the essence of absolute reliability.\newline\null\hfill\href{https://www.valmiki.iitk.ac.in/sloka?field_kanda_tid=6&language=dv&field_sarga_value=74}{\textit{Yuddha Kanda, 74:71-74}}
\end{tcolorbox}
\end{minipage}

\vspace{0.5em}
\subsubsection*{\textbf{Visual Rhythm Map:}}
\begin{table}[h!]
\centering
\renewcommand{\arraystretch}{1.2}
\setlength{\tabcolsep}{0pt}
\begin{tabularx}{\textwidth}{|*{16}{>{\centering\arraybackslash\hsize=1.0\hsize}X|}}
\hline
\multicolumn{3}{|>{\centering\arraybackslash\hsize=3.0\hsize}X|}{\textbf{\guru{ला}\laghu{य}} \par (3)} &
\multicolumn{5}{>{\centering\arraybackslash\hsize=5.0\hsize}X|}{\textbf{\laghu{स}\guru{जी}\laghu{व}\laghu{न}} \par (5)} &
\multicolumn{3}{>{\centering\arraybackslash\hsize=3.0\hsize}X|}{\textbf{\laghu{ल}\laghu{ख}\laghu{न}} \par (3)} &
\multicolumn{5}{>{\centering\arraybackslash\hsize=5.0\hsize}X|}{\textbf{\laghu{जि}\guru{या}\guru{ये}} \par (5)} \\
\hline
\end{tabularx}
\vspace{-1.5em}
\end{table}

\begin{table}[h!]
\centering
\renewcommand{\arraystretch}{1.2}
\setlength{\tabcolsep}{0pt}
\begin{tabularx}{\textwidth}{|*{16}{>{\centering\arraybackslash\hsize=1.0\hsize}X|}}
\hline
\multicolumn{7}{|>{\centering\arraybackslash\hsize=7.0\hsize}X|}{\textbf{\guru{श्री}\laghu{र}\laghu{घु}\guru{बी}\laghu{र}} \par (7)} &
\multicolumn{3}{>{\centering\arraybackslash\hsize=3.0\hsize}X|}{\textbf{\laghu{ह}\laghu{र}\laghu{षि}} \par (3)} &
\multicolumn{2}{>{\centering\arraybackslash\hsize=2.0\hsize}X|}{\textbf{\laghu{उ}\laghu{र}} \par (2)} &
\multicolumn{4}{>{\centering\arraybackslash\hsize=4.0\hsize}X|}{\textbf{\guru{ला}\guru{ये}} \par (4)} \\
\hline
\end{tabularx}
\vspace{-1.5em}
\end{table}

\vspace{1em}
\textbf{ \deva{सरल-संस्कृतम्}:}\\
 \deva{भवान् सञ्जीवनी-बूटीम् आनीय लक्ष्मणस्य प्राणान् रक्षितवान् (जियाये)।}\\
 \deva{तत् दृष्ट्वा श्रीरामः (रघुवीरः) हर्षितः सन् भवन्तं वक्षसि (उर) स्थापितवान् (आलिङ्गनं कृतवान्)॥}\\


Bringing the Sanjeevani herb, you resurrected Lakshmana.\\
Seeing this, Lord Rama joyfully embraced you to his chest/heart.


\vspace{1em}
\newpage
\textbf{\deva{प्रतिपदार्थः} (Meaning of each word):}
\begin{table}[h!]
\centering
\renewcommand{\arraystretch}{1.2}
\begin{tabularx}{\textwidth}{@{} l l X X @{}}
\toprule
\textbf{Awadhi} & \textbf{Sanskrit} & \textbf{English} & \textbf{Grammar \& Notes} \\
\midrule
\deva{लाय} / \telu{లాయ} / \textit{lāya}  &  \deva{आनीय}  &  Having brought  &   \\
\deva{सजीवन} \gotchaicon / \telu{సజీవన} / \textit{sajīvana}  &  \deva{सञ्जीवनीम्}  &  Sanjeevani herb  &   \\
\deva{लखन} / \telu{లఖన} / \textit{lakhana}  &  \deva{लक्ष्मणम्}  &  Lakshmana  & \\ 
\deva{जियाये} / \telu{జియాయే} / \textit{jiyāye}  &  \deva{जीवितवान्}  &  Resurrected / Revived  &   \\
\deva{श्रीरघुबीर} / \telu{శ్రీరఘుబీర} / \textit{śrīraghubīra}  &  \deva{श्रीरघुवीरः}  &  Lord Rama  & \\ 
\deva{हरषि} / \telu{హరషి} / \textit{haraṣi}  &  \deva{हर्षितः}  &  Joyfully  & Short ending 'i' \\
\deva{उर} / \telu{ఉర} / \textit{ura}  &  \deva{उरसि / वक्षसि}  &  In the chest/heart  &   \\
\deva{लाये} / \telu{లాయే} / \textit{lāye}  &  \deva{स्थापितवान्}  &  Brought to (embraced)  &   \\
\bottomrule
\end{tabularx}
\end{table}

\begin{tcolorbox}[gotchabox]
\begin{itemize}
    \item \textbf{Phonetic Shifts:} \deva{रघुवीर} (Raghuveera) hardens into \deva{रघुबीर} (Raghubeera). The Sanskrit `jñ/ñj` cluster in \deva{सञ्जीवनी} (Sanjeevani) simplifies into the pure Awadhi \deva{सजीवन} (Sajeevan).
\end{itemize}
\end{tcolorbox}



\newpage
